% from ega4-14-fr.tex
\oldpage[IV-3]{206}

De plus, on peut, dans la condition précédente, se borner au cas où $A$ est
complet, a un corps résiduel algébriquement clos, et où $x'$ est rationnel sur
$k(y')$.
\end{proposition}

\begin{proof}
Si $Y'$ est comme dans l'énoncé, la nécessité de la condition résulte de ce que
$f'$ doit être ouvert au point $x'$, et du critère (1.10.3). Pour voir que la
condition est suffisante, considérons un morphisme de type fini $g:Y''\to U$, et
soient $X''=X\times_Y Y''$, $f''=f_{(Y'')}:X''\to Y''$, et $x''$ un point de $X''$
au-dessus de $x$. Posons $y''=f''(x'')$, et soit $t$ une générisation de $y''$
dans $Y''$, distincte de $y''$. Puisque $Y''$ est localement noethérien, il
résulte de (\textbf{II}, 7.1.9) et ($\mathbf{0}_{\mathrm{III}}$, 10.3.1) qu'il
existe un schéma $Y'=\operatorname{Spec}(A)$, où $A$ est un anneau de valuation
discrète complet, dont le corps résiduel est une clôture algébrique de $k(x'')$,
et un morphisme $g:Y'\to Y''$ tels que, si $s$ et $y'$ sont le point générique et
le point fermé de $Y'$, on ait $g(s)=t$ et $g(y')=y''$. Il y a alors un point $x'$
de $X'=X\times_Y Y'=X''\times_{Y''}Y'$ dont les projections dans $X''$ et $Y'$
soient $x''$ et $y'$ respectivement, et qui soit rationnel sur $k(y')$ (I, 3.4.9).
L'hypothèse entraîne qu'il y a une générisation $z'$ de $x'$ dans $X'$ dont la
projection dans $Y'$ soit $s$~; si $z''$ est la projection de $z'$ dans $X''$,
$z''$ est une générisation de $x''$ et sa projection dans $Y''$ est $t$~; on
conclut donc de (1.10.3) que $f''$ est ouvert au point $x''$, donc que $f$ est
universellement ouvert au point $x$ (14.3.3.1, (i)).
\end{proof}

\begin{corollary}[14.3.7]
\label{IV.14.3.7-fr}
Les notations étant celles de (14.3.6)~:
\begin{enumerate}[label=(\roman*)]
  \item Étant donné un point $y\in Y$, pour que $f$ soit universellement ouvert en
  tous les points de $f^{-1}(y)$, il faut et il suffit que pour tout morphisme
  $g:Y'\to Y$, où $Y'$ est le spectre d'un anneau de valuation discrète, et où
  l'image $g(y')$ du point fermé $y'$ de $Y'$ est $y$, toute composante
  irréductible de $X'=X\times_Y Y'$ domine $Y'$.
  \item Pour que $f$ soit universellement ouvert, il faut et il suffit que pour
  tout morphisme $g:Y'\to Y$, où $Y'$ est le spectre d'un anneau de valuation
  discrète, toute composante irréductible de $X'=X\times_Y Y'$ domine $Y'$.
\end{enumerate}
\end{corollary}

\begin{proof}
Il est clair qu'il suffit de prouver (i)~; la nécessité de (i) résulte de (14.3.2)
et sa suffisance de (14.3.6).
\end{proof}

\begin{proposition}[14.3.8]
\EGAReaderTarget{ega4.c04.prop.14-3-8}{IV.14.3.8-fr}
Soient $Y$ un préschéma localement noethérien, irréductible, régulier et de
dimension 1 (par exemple le spectre d'un anneau de Dedekind), $f:X\to Y$ un
morphisme localement de type fini, $y$ un point de $Y$. Les conditions suivantes
sont équivalentes~:
\begin{enumerate}[label=\emph{\alph*)}]
  \item $f_{\mathrm{red}}$ est plat en tout point de $f^{-1}(y)$.
  \item $f$ est universellement ouvert dans un voisinage de $f^{-1}(y)$.
  \item $f$ est ouvert dans un voisinage de $f^{-1}(y)$.
  \item Toute composante irréductible de $X$ rencontrant $f^{-1}(y)$ domine $Y$.
\end{enumerate}
\end{proposition}

\begin{proof}
Comme $f_{\mathrm{red}}$ est localement de type fini (1.3.4), $a)$ entraîne que
$f_{\mathrm{red}}$ est plat dans un voisinage de $f^{-1}(y)$ (11.1.1), et il
suffit d'appliquer (2.4.6) dans un tel voisinage pour voir que $a)$ entraîne $b)$.
L'implication $b)\Rightarrow c)$ est triviale, et l'implication
$c)\Rightarrow d)$ résulte de (1.10.4) appliqué à un voisinage de $f^{-1}(y)$.
Reste à voir que $d)$ entraîne $a)$. On peut évidemment, en vertu de (1.3.4), se
borner au cas où $X$ est réduit. La question étant en outre locale sur $X$ et sur
$Y$, on peut supposer $Y=\operatorname{Spec}(A)$ et
$X=\operatorname{Spec}(B)$
